% Generated by `taskgen course examples` from dlp.shanks (seed 3).
% Do not edit: change examples.yaml and regenerate.
\begin{example}[алгоритм Шенкса]\label{ex:bsgs-shanks}
Решить задачу дискретного логарифмирования методом Шенкса: \[8^x \equiv 27\ (\textrm{mod}\ 71).\] Обосновать выбор длины списков.

\par\noindent\textit{Решение.} Порядок $g$ делит $p - 1 = 70$; при длине списков $m = \lceil\sqrt{70}\rceil = 9$ любой показатель $x < m^2$ записывается как $x = im + j$, $0 \le i, j < m$ (а $m^2 \ge p - 1$). Малые шаги $g^j$:\[\begin{array}{c|rrrrrrrrr}j & 0 & 1 & 2 & 3 & 4 & 5 & 6 & 7 & 8 \\ \hline g^j & 1 & 8 & 64 & 15 & 49 & 37 & 12 & 25 & 58\end{array}\]$g^{-1} \equiv 9$, $g^{-m} \equiv 9^{9} \equiv 43$. Большие шаги $h \cdot g^{-im}$ до совпадения:\[\begin{array}{c|rr}i & 0 & 1 \\ \hline hg^{-im} & 27 & 25\end{array}\]Совпадение $h g^{-1 \cdot 9} = g^{7}$, откуда $x = 1 \cdot 9 + 7 = 16$. $\ord(8) = 35$: $8^{35} \equiv 1$, $8^{7} \equiv 25$, $8^{5} \equiv 37$.\[x = 16\ (\textrm{mod}\ 35).\]
\end{example}
